\section{A grazing-uniform curvature graph transform}
\label{sec:uc-curvature}
The preceding numbered proof bodies are retained verbatim. We now discharge
one of their explicitly exposed premises: admissibility of the actual angular
source for homogeneous curve growth. This is not a deduction of roofwise
pressure from a curvature bound. We prove curvature invariance here,
one-step weighted complexity and distortion in the next section, and only
then apply a precisely identified curve-sum lemma. No event or collision is
changed by these constructions.

Write $t_*=\tau_-=3/50$, $k_-=100/47$, $k_+=20/9$, and
$U=k_++t_*^{-1}=170/9$. At a fixed radius write $\kappa=1/R$.
Our curves are graphs $\varphi=\varphi(s)$ in actual boundary arclength,
with $V=d\varphi/ds\in[\kappa,U]$. Their signed geometric curvature is
\[
              K=\frac{dV/ds}{(1+V^2)^{3/2}}.
\]
Neither $K$ nor $dV/ds$ is an optical wavefront curvature.
The notation in this section is local to one actual regular collision.

\begin{lemma}[Exact circular graph transform]\label{lem:uc-transform}
Let $c=\cos\varphi$, $p=\sin\varphi$, $d=\cos\varphi_1$,
$p_1=\sin\varphi_1$ at the two actual collisions, and let $\tau$ be their
physical flight time. Put
\begin{equation}\label{eq:uc-wa}
 w=\kappa+V,\qquad A=c+\tau w,\qquad y=dw/A.
\end{equation}
On a regular branch, $c,d>0$, and the image slope and signed derivatives are
\begin{align}
 s_{1,s}&=-A/d,& V_1&=\kappa+y, &
 \tau_s&=-p-p_1 A/d,\label{eq:uc-first}\\
 V_{1,s_1}
 &=-\frac{p_1w}{A}(\kappa+2y)
   -\frac{pd^2w(\kappa+2V)}{A^3}
   -\frac{cd^2}{A^3}V_s.\label{eq:uc-second}
\end{align}
In particular, $\kappa\le V_1\le U$.
\end{lemma}
\begin{proof}
Let $n,t$ and $n_1,t_1$ be the outward normal and its counterclockwise
rotation at the two contacts. The outgoing velocity is $cn+pt$, and
at the next disk its incoming representation is $-dn_1+p_1t_1$.
The outgoing direction derivative is $\vartheta_s=\kappa+V=w$.
Differentiate $q_1=q+\tau v$ and project perpendicular to $v$.
The initial projection is $c$, the terminal one is $-d s_{1,s}$,
so $-d s_{1,s}=c+\tau w$. Projection along $v$ gives
$\tau_s=p_1s_{1,s}-p$. Specular reflection gives
$d\varphi_1=\kappa\,ds_1-d\vartheta$, whence
$V_1=\kappa-w/s_{1,s}$. This proves \eqref{eq:uc-first}.

For the second derivative use
\[
 c_s=-pV,\quad d_s=p_1V_1A/d,\quad
 A_s=-p(\kappa+2V)-p_1Aw/d+\tau V_s.
\]
Differentiate $V_1=\kappa+dw/A$ and divide by $s_{1,s}=-A/d$.
The coefficient of $V_s$ is
$-d^2/A^2+d^2\tau w/A^3=-cd^2/A^3$; the remaining terms give
\eqref{eq:uc-second}. Finally $0<y\le1/\tau\le1/t_*$,
so the image stays in the stated cone. Differentiation is only on
regular branches; no value of $DT$ at a tangency is asserted.
\end{proof}

\begin{theorem}[Uniform contraction of geometric curvature]
\label{thm:uc-curvature}
For every radius in the original interval and every regular actual flight,
\begin{equation}\label{eq:uc-affine}
 |K_1|\le q_*|K|+C_*,\qquad
 q_*:=\frac{4\bigl(U+(2U)^{-1}\bigr)^3}
 {27k_-(k_-+U)^2\{1+t_*(k_-+U)\}}<\frac12,
 \quad C_*:=\frac{58}{27t_*k_-^2}.
\end{equation}
These constants are independent of both positive incidences and of the
flight count. In exact rational form,
\[
 q_*={946236091513673179721\over2005319667933288000000},
 \qquad C_*={64061\over8100}.
\]
Consequently any
\begin{equation}\label{eq:uc-B}
       B_{\rm cv}\ge\max\{K_0,2C_*\}
\end{equation}
defines a forward invariant class of homogeneous $C^2$ unstable graphs
of geometric curvature at most $B_{\rm cv}$. It contains the actual sources of
Theorem~\ref{thm:ag-family}. Along every regular branch of length $l$,
\begin{equation}\label{eq:uc-iterate}
 |K_l|\le q_*^l|K_0^{\rm curve}|+
              C_*\frac{1-q_*^l}{1-q_*}.
\end{equation}
This assertion includes arbitrarily close approaches to grazing but not a
smooth extension through a singular collision.
\end{theorem}
\begin{proof}
After division of \eqref{eq:uc-second} by $(1+V_1^2)^{3/2}$, the
absolute coefficient of $K$ is
\begin{align}
 Q&=\frac{cd^2(1+V^2)^{3/2}}
       {\{A^2+(\kappa A+dw)^2\}^{3/2}}\notag\\
  &=\frac cA\frac{(1+V^2)^{3/2}}{w^2}
                  \frac{y^2}{\{1+(\kappa+y)^2\}^{3/2}}.
                                                    \label{eq:uc-Q}
\end{align}
Here $c/A\le[1+t_*w]^{-1}$ and
\[
 \frac{y^2}{\{1+(\kappa+y)^2\}^{3/2}}
 \le\frac{y^2}{(\kappa+y)^3}\le\frac4{27\kappa}.
\]
The last inequality follows by differentiating $y^2/(\kappa+y)^3$;
its maximum for $y\ge0$ is at $y=2\kappa$.
For fixed $\kappa$, the function
\[
 h(V)=\frac{(1+V^2)^{3/2}}
 {(\kappa+V)^2\{1+t_*(\kappa+V)\}}
\]
is increasing for $V\ge\kappa>1$. Indeed
$3V/(1+V^2)>3/(\kappa+V)$ because $\kappa V>1$, and the
remaining lower bound for its logarithmic derivative is
$1/(\kappa+V)-t_*/[1+t_*(\kappa+V)]>0$.
For fixed $V$ it decreases with $\kappa$. Thus
$Q\le 4h(U)|_{\kappa=k_-}/(27k_-)$.
The elementary bound $\sqrt{1+U^2}\le U+(2U)^{-1}$ gives $Q\le q_*$.
The displayed rational evaluation verifies $q_*<1/2$ without rounding.

For the first forcing term in \eqref{eq:uc-second}, use $w/A\le1/t_*$
and $\kappa+2y\le2(\kappa+y)$ to bound its normalized magnitude by
$2/(t_*k_-^2)$. For the second, write its normalized magnitude as
\[
 \frac{\kappa+2V}{Aw}
            \frac{y^2}{\{1+(\kappa+y)^2\}^{3/2}}
 \le\frac{2}{t_*w}\frac4{27\kappa}
 \le\frac4{27t_*k_-^2}.
\]
Their sum is $C_*$. This proves \eqref{eq:uc-affine}. Induction gives
\eqref{eq:uc-iterate}, and $q_*B_{\rm cv}+C_*\le B_{\rm cv}$ follows from
\eqref{eq:uc-B}. A homogeneous cut or length subdivision restricts a
smooth curve and does not increase its curvature. The exact source
curvature was bounded by $K_0$ in Theorem~\ref{thm:ag-density}.
All denominators used here are positive on the regular domain; the
uniform estimates themselves require no incidence cutoff.
\end{proof}

\begin{remark}[What the norm conversion does and does not do]
The graph second derivative on this invariant class is at most
$B_2=B_{\rm cv}(1+U^2)^{3/2}$. This is a consequence of geometric curvature
invariance, not a claim that the same bound on $\varphi''$ is itself
an invariant rectangular constraint. The class in (H2) of
\cite{DZUniform} is formulated with geometric curvature; equivalence
of these two boundedness notions is used only for derivative estimates.
The optical Jacobi curvature and the long-record roof Hessian are different
quantities and are not controlled by \eqref{eq:uc-affine}.
\end{remark}
